\unitchapter{2}{9}{Data Communication \& Computer Networks}{p2-09}

\unitmeta{Expected questions: 10--12 · Target: 10+ · Type: subnetting, sliding
window, Shannon/Nyquist numericals + Forouzan concepts}

\begin{sylbox}

\begin{itemize}
\tightlist
\item[$\square$]
  Data communication: components, data flow, topologies, protocols \& standards, OSI \& TCP/IP
\item[$\square$]
  Analog \& digital signals: bandwidth, impairments, Nyquist \& Shannon limits, performance
\item[$\square$]
  Digital transmission: line coding, block coding, scrambling, PCM, delta modulation
\item[$\square$]
  Analog transmission: ASK, FSK, PSK, QAM
\item[$\square$]
  Multiplexing (FDM, WDM, TDM), spread spectrum, transmission media
\item[$\square$]
  Switching: circuit, message, packet (datagram, virtual circuit)
\item[$\square$]
  Data link layer: framing, error detection \& correction, flow control, Stop-and-Wait, Go-Back-N, Selective Repeat, HDLC, PPP
\item[$\square$]
  MAC: ALOHA, CSMA, CSMA/CD, CSMA/CA, polling, token passing, FDMA/TDMA/CDMA; Ethernet
\item[$\square$]
  Network layer: IPv4, subnetting, CIDR, NAT, IPv6, routing (DV, LS, path vector), RIP, OSPF, BGP, ARP, ICMP, DHCP
\item[$\square$]
  Transport layer: UDP, TCP, SCTP, congestion control, QoS
\item[$\square$]
  Application layer: DNS, HTTP, FTP, SMTP, TELNET, SNMP
\item[$\square$]
  Network security: cryptography, RSA, Diffie-Hellman, digital signatures, VPN, firewalls
\item[$\square$]
  Mobile technology: GSM, CDMA, GPRS, Mobile IP, MANETs, satellites, wireless LANs
\item[$\square$]
  Cloud computing \& IoT
\end{itemize}

\end{sylbox}

\needspace{246pt}

\hypertarget{p2-09--1-network-models}{%
\section{Network Models}\label{p2-09--1-network-models}}

\needspace{210pt}

\hypertarget{p2-09--11-osi-vs-tcpip}{%
\subsection{OSI vs TCP/IP}\label{p2-09--11-osi-vs-tcpip}}

\begin{listingblock}{160pt}
\begin{Verbatim}[fontsize=\fontsize{6.8}{8.2}\selectfont]
   OSI (7 layers)              TCP/IP (4/5 layers)        PDU          Devices / protocols
 ┌──────────────────┐        ┌───────────────────┐
 │ 7 Application    │        │                   │
 ├──────────────────┤        │                   │                     HTTP, FTP, SMTP, DNS,
 │ 6 Presentation   │  ───►  │   Application     │      Message        SNMP, TELNET, DHCP
 ├──────────────────┤        │                   │                     (gateway)
 │ 5 Session        │        │                   │
 ├──────────────────┤        ├───────────────────┤
 │ 4 Transport      │  ───►  │   Transport       │      Segment        TCP, UDP, SCTP
 ├──────────────────┤        ├───────────────────┤      (datagram)
 │ 3 Network        │  ───►  │   Internet        │      Packet         IP, ICMP, ARP, IGMP (router)
 ├──────────────────┤        ├───────────────────┤
 │ 2 Data Link      │  ───►  │   Network access  │      Frame          Ethernet, PPP, HDLC (switch, bridge)
 ├──────────────────┤        │   (link+physical) │
 │ 1 Physical       │        │                   │      Bits           (hub, repeater)
 └──────────────────┘        └───────────────────┘
 Mnemonic (top→bottom): "All People Seem To Need Data Processing"
\end{Verbatim}
\end{listingblock}

\begingroup\small

\begin{longtable}[]{@{}
  >{\raggedright\arraybackslash}p{(\columnwidth - 4\tabcolsep) * \real{0.1270}}
  >{\raggedright\arraybackslash}p{(\columnwidth - 4\tabcolsep) * \real{0.6448}}
  >{\raggedright\arraybackslash}p{(\columnwidth - 4\tabcolsep) * \real{0.2282}}@{}}
\toprule\noalign{}
\begin{minipage}[b]{\linewidth}\raggedright
\textbf{Layer}
\end{minipage} & \begin{minipage}[b]{\linewidth}\raggedright
\textbf{Key responsibilities}
\end{minipage} & \begin{minipage}[b]{\linewidth}\raggedright
\textbf{Address}
\end{minipage} \\
\midrule\noalign{}
\endhead
\bottomrule\noalign{}
\endlastfoot
Physical & Bit transmission, encoding, topology, data rate & --- \\
Data link & \textbf{Framing}, physical addressing, flow \& error control (hop-to-hop), MAC & \textbf{MAC (48-bit)} \\
Network & \textbf{Logical addressing, routing} (host-to-host / source-to-destination) & \textbf{IP} \\
Transport & \textbf{Process-to-process (end-to-end)} delivery, segmentation, flow, error, congestion control & \textbf{Port (16-bit)} \\
Session & Dialog control, \textbf{synchronisation (checkpoints)} & --- \\
Presentation & \textbf{Translation, encryption, compression} & --- \\
Application & User services & Specific (URL, email) \\
\end{longtable}

\endgroup

\needspace{130pt}

\hypertarget{p2-09--12-topologies}{%
\subsection{Topologies}\label{p2-09--12-topologies}}

\begingroup\small

\begin{longtable}[]{@{}
  >{\raggedright\arraybackslash}p{(\columnwidth - 4\tabcolsep) * \real{0.1385}}
  >{\raggedright\arraybackslash}p{(\columnwidth - 4\tabcolsep) * \real{0.3994}}
  >{\raggedright\arraybackslash}p{(\columnwidth - 4\tabcolsep) * \real{0.3448}}@{}}
\toprule\noalign{}
\begin{minipage}[b]{\linewidth}\raggedright
\textbf{Topology}
\end{minipage} & \begin{minipage}[b]{\linewidth}\raggedright
\textbf{Cables / links (n devices)}
\end{minipage} & \begin{minipage}[b]{\linewidth}\raggedright
\textbf{Notes}
\end{minipage} \\
\midrule\noalign{}
\endhead
\bottomrule\noalign{}
\endlastfoot
\textbf{Mesh} & \textbf{n(n−1)/2} links; each device n−1 I/O ports & Robust, expensive \\
Star & n links & Central hub; hub failure kills network \\
Bus & 1 backbone + drop lines & Cheap; backbone fault fails all \\
Ring & n links & Token passing; one break can disable \\
Tree / Hybrid & --- & \\
\end{longtable}

\endgroup

Data flow: \textbf{simplex} (keyboard), \textbf{half-duplex} (walkie-talkie), \textbf{full-duplex} (telephone).

\needspace{136pt}

\hypertarget{p2-09--2-signals--data-rates}{%
\section{Signals \& Data Rates}\label{p2-09--2-signals--data-rates}}

\begin{listingblock}{86pt}
\begin{Verbatim}[fontsize=\fontsize{7.7}{9.3}\selectfont]
Nyquist (noiseless):    Max bit rate = 2 · B · log₂ L          (L signal levels)
Shannon (noisy):        Capacity C = B · log₂(1 + SNR)
SNR_dB = 10 log₁₀(SNR)  → 30 dB = 1000, 20 dB = 100, 10 dB = 10
Bandwidth-delay product = bandwidth × propagation delay  (bits "in the pipe")
Propagation delay = distance / propagation speed ;  Transmission delay = frame size / bandwidth
Throughput, latency = propagation + transmission + queuing + processing
Baud rate (signal rate) S = N / r  where r = bits per signal element
\end{Verbatim}
\end{listingblock}

\textbf{Worked}: telephone line B = 3000 Hz, SNR = 3162 (≈ 35 dB): C = 3000 × log₂(3163) ≈ 3000 × 11.62 = \textbf{34.86 kbps}.
Nyquist: B = 3000 Hz, 4 levels → 2 × 3000 × 2 = \textbf{12 kbps}.

Impairments: \textbf{attenuation} (loss of energy, dB = 10 log(P₂/P₁)), \textbf{distortion} (different delays of components), \textbf{noise} (thermal, induced, crosstalk, impulse).

\needspace{216pt}

\hypertarget{p2-09--3-digital--analog-transmission}{%
\section{Digital \& Analog Transmission}\label{p2-09--3-digital--analog-transmission}}

\needspace{180pt}

\hypertarget{p2-09--31-line-coding}{%
\subsection{Line Coding}\label{p2-09--31-line-coding}}

\begin{listingblock}{130pt}
\begin{Verbatim}[fontsize=\fontsize{7.5}{9.1}\selectfont]
 Bits:        1      0      1      1      0
 NRZ-L      ──────┐      ┌─────────────┐      
 (1 = high)       └──────┘             └──────

 Bits:        1       0       1       1       0
 Manchester    ┌───────┐       ┌───┐   ┌───────┐   
 (IEEE)     ───┘       └───────┘   └───┘       └───
 0 = high→low, 1 = low→high: every bit has a MID-BIT transition → receiver recovers the clock

 NRZ-I: invert the level at the start of every 1, no change for 0
 Differential Manchester: transition at START of bit for 0, none for 1; mid-bit transition always
 AMI (bipolar): 0 → zero voltage; successive 1s alternate +V and −V (no DC component)
\end{Verbatim}
\end{listingblock}

\begingroup\small

\begin{longtable}[]{@{}
  >{\raggedright\arraybackslash}p{(\columnwidth - 4\tabcolsep) * \real{0.3333}}
  >{\raggedright\arraybackslash}p{(\columnwidth - 4\tabcolsep) * \real{0.1477}}
  >{\raggedright\arraybackslash}p{(\columnwidth - 4\tabcolsep) * \real{0.5190}}@{}}
\toprule\noalign{}
\begin{minipage}[b]{\linewidth}\raggedright
\textbf{Scheme}
\end{minipage} & \begin{minipage}[b]{\linewidth}\raggedright
\textbf{Category}
\end{minipage} & \begin{minipage}[b]{\linewidth}\raggedright
\textbf{Notes}
\end{minipage} \\
\midrule\noalign{}
\endhead
\bottomrule\noalign{}
\endlastfoot
NRZ-L, NRZ-I & Polar & No self-sync; baseline wandering \\
RZ & Polar & 3 levels; 2 transitions per bit \\
\textbf{Manchester}, Differential Manchester & Biphase & Self-synchronising; bandwidth = 2× NRZ; Ethernet (10 Mbps) / Token ring \\
AMI, Pseudoternary & Bipolar & No DC component \\
2B1Q, 8B6T, 4D-PAM5 & Multilevel & \\
MLT-3 & Multitransition & 100Base-TX \\
\end{longtable}

\endgroup

\begin{itemize}
\tightlist
\item
  \textbf{Block coding}: 4B/5B (Fast Ethernet), 8B/10B (Gigabit) --- for synchronisation \& error detection.
\item
  \textbf{Scrambling}: B8ZS (North America), HDB3 (Europe) --- replace long zeros in AMI.
\item
  \textbf{PCM}: Sampling → Quantisation → Encoding. \textbf{Nyquist sampling theorem}: fs ≥ 2·f\_max. Bit rate = fs × n bits. Voice: 8000 samples/s × 8 bits = \textbf{64 kbps}.
\item
  \textbf{Delta modulation}: 1 bit per sample (up/down).
\item
  Transmission modes: \textbf{parallel}; \textbf{serial} --- asynchronous (start/stop bits, gaps), synchronous, isochronous.
\end{itemize}

\needspace{95pt}

\hypertarget{p2-09--32-digital-to-analog}{%
\subsection{Digital-to-Analog}\label{p2-09--32-digital-to-analog}}

\diagram{figures/p2-09-00}

\begin{itemize}
\tightlist
\item
  QPSK: 2 bits per symbol; 16-QAM: 4 bits per symbol; 64-QAM: 6. Bits per baud = log₂(constellation points).
\item
  Analog-to-analog: AM, FM, PM. FM bandwidth (Carson) = 2(1 + β)B.
\end{itemize}

\needspace{130pt}

\hypertarget{p2-09--4-multiplexing--media}{%
\section{Multiplexing \& Media}\label{p2-09--4-multiplexing--media}}

\begingroup\small

\begin{longtable}[]{@{}
  >{\raggedright\arraybackslash}p{(\columnwidth - 4\tabcolsep) * \real{0.1759}}
  >{\raggedright\arraybackslash}p{(\columnwidth - 4\tabcolsep) * \real{0.1960}}
  >{\raggedright\arraybackslash}p{(\columnwidth - 4\tabcolsep) * \real{0.5787}}@{}}
\toprule\noalign{}
\begin{minipage}[b]{\linewidth}\raggedright
\textbf{Technique}
\end{minipage} & \begin{minipage}[b]{\linewidth}\raggedright
\textbf{Domain}
\end{minipage} & \begin{minipage}[b]{\linewidth}\raggedright
\textbf{Use}
\end{minipage} \\
\midrule\noalign{}
\endhead
\bottomrule\noalign{}
\endlastfoot
FDM & Frequency (analog) & Radio, TV, cable \\
WDM & Wavelength (optical) & Fibre (DWDM) \\
Synchronous TDM & Time slots fixed & T1 (24 channels, 1.544 Mbps), E1 (32 channels, 2.048 Mbps) \\
Statistical TDM & Slots on demand & Needs address in slot \\
Spread spectrum & FHSS, DSSS & Wireless, security \\
\end{longtable}

\endgroup

\begin{itemize}
\tightlist
\item
  T1 = (24 × 8 + 1 framing bit) × 8000 = \textbf{1.544 Mbps}.
\end{itemize}

\textbf{Transmission media}:

\begin{listingblock}{71pt}
\begin{Verbatim}[fontsize=\fontsize{7.0}{8.5}\selectfont]
Guided:   Twisted pair (UTP Cat5/6, RJ-45; twisting reduces crosstalk)
          Coaxial (BNC; cable TV)
          Optical fibre (total internal reflection; single-mode vs multimode; highest bandwidth, no EMI)
Unguided: Radio waves (3 kHz–1 GHz, omnidirectional)
          Microwaves (1–300 GHz, unidirectional, line-of-sight)
          Infrared (300 GHz–400 THz, short range, can't penetrate walls)
\end{Verbatim}
\end{listingblock}

\needspace{130pt}

\hypertarget{p2-09--5-switching}{%
\section{Switching}\label{p2-09--5-switching}}

\begingroup\small

\begin{longtable}[]{@{}
  >{\raggedright\arraybackslash}p{(\columnwidth - 6\tabcolsep) * \real{0.2166}}
  >{\raggedright\arraybackslash}p{(\columnwidth - 6\tabcolsep) * \real{0.2861}}
  >{\raggedright\arraybackslash}p{(\columnwidth - 6\tabcolsep) * \real{0.2634}}
  >{\raggedright\arraybackslash}p{(\columnwidth - 6\tabcolsep) * \real{0.2338}}@{}}
\toprule\noalign{}
\begin{minipage}[b]{\linewidth}\raggedright
\textbf{Circuit switching}
\end{minipage} & \begin{minipage}[b]{\linewidth}\raggedright
\textbf{Packet (datagram)}
\end{minipage} & \begin{minipage}[b]{\linewidth}\raggedright
\textbf{Virtual circuit}
\end{minipage} & \begin{minipage}[b]{\linewidth}\raggedright
\textbf{Message switching}
\end{minipage} \\
\midrule\noalign{}
\endhead
\bottomrule\noalign{}
\endlastfoot
Dedicated path, setup/\allowbreak{}teardown & No setup; each packet routed independently & Setup phase; packets follow same path & Store-and-forward whole message \\
Telephone & Internet (IP) & ATM, Frame Relay, X.25 & Telegraph \\
In-order, constant delay & Out-of-order possible & In-order & Large buffers \\
Physical layer & Network layer & Data link (ATM) / network & Application \\
\end{longtable}

\endgroup

\textbf{Delay for packet switching}: message split into k packets over h hops: total ≈ (k + h − 1) × transmission time per packet (+ propagation).

\needspace{197pt}

\hypertarget{p2-09--6-data-link-layer}{%
\section{Data Link Layer}\label{p2-09--6-data-link-layer}}

\needspace{161pt}

\hypertarget{p2-09--61-framing}{%
\subsection{Framing}\label{p2-09--61-framing}}

Character count, \textbf{byte stuffing} (ESC/FLAG), \textbf{bit stuffing} (insert 0 after five consecutive 1s --- HDLC flag 01111110), physical layer coding violations.
Bit-stuffing example: data \texttt{0111111111110} → \texttt{011111011111010}.

\needspace{95pt}

\hypertarget{p2-09--62-error-detection--correction}{%
\subsection{Error Detection \& Correction}\label{p2-09--62-error-detection--correction}}

\begin{itemize}
\tightlist
\item
  \textbf{Parity}: detects odd number of bit errors. 2-D parity detects all 1, 2, 3-bit errors and corrects single.
\item
  \textbf{Checksum} (1\textquotesingle s complement sum; used in IP, TCP, UDP).
\item
  \textbf{CRC}: append (degree r) zeros, divide by generator (mod-2), remainder = CRC. Detects all burst errors of length ≤ r; all odd errors if generator has factor (x + 1).
\end{itemize}

\textbf{CRC worked}: data 1101011011, generator x⁴ + x + 1 = 10011 → append 4 zeros, divide → remainder \textbf{1110}; transmitted \textbf{11010110111110}.

\begin{itemize}
\tightlist
\item
  \textbf{Hamming distance}: detect d errors needs d\_min ≥ d + 1; correct t errors needs d\_min ≥ 2t + 1.
\item
  Hamming code (7,4): r parity bits with 2ʳ ≥ m + r + 1.
\end{itemize}

\needspace{147pt}

\hypertarget{p2-09--63-flow--error-control-sliding-window}{%
\subsection{Flow \& Error Control (Sliding Window)}\label{p2-09--63-flow--error-control-sliding-window}}

\begin{listingblock}{97pt}
\begin{Verbatim}[fontsize=\fontsize{9.0}{10.9}\selectfont]
a = propagation delay / transmission time = Tp / Tt
Stop-and-Wait efficiency   η = 1 / (1 + 2a)
Sliding window (size W)    η = W / (1 + 2a)   if W < 1 + 2a, else 1
For 100% utilisation       W ≥ 1 + 2a
Sequence number bits n:    Go-Back-N  W_s ≤ 2ⁿ − 1, W_r = 1
                           Selective Repeat  W_s = W_r ≤ 2ⁿ⁻¹
Throughput = η × bandwidth
\end{Verbatim}
\end{listingblock}

\diagram{figures/p2-09-01}

\begingroup\small

\begin{longtable}[]{@{}
  >{\raggedright\arraybackslash}p{(\columnwidth - 8\tabcolsep) * \real{0.1753}}
  >{\raggedright\arraybackslash}p{(\columnwidth - 8\tabcolsep) * \real{0.1783}}
  >{\raggedright\arraybackslash}p{(\columnwidth - 8\tabcolsep) * \real{0.1982}}
  >{\raggedright\arraybackslash}p{(\columnwidth - 8\tabcolsep) * \real{0.2151}}
  >{\raggedright\arraybackslash}p{(\columnwidth - 8\tabcolsep) * \real{0.2331}}@{}}
\toprule\noalign{}
\begin{minipage}[b]{\linewidth}\raggedright
\textbf{Protocol}
\end{minipage} & \begin{minipage}[b]{\linewidth}\raggedright
\textbf{Sender window}
\end{minipage} & \begin{minipage}[b]{\linewidth}\raggedright
\textbf{Receiver window}
\end{minipage} & \begin{minipage}[b]{\linewidth}\raggedright
\textbf{Retransmits}
\end{minipage} & \begin{minipage}[b]{\linewidth}\raggedright
\textbf{ACK}
\end{minipage} \\
\midrule\noalign{}
\endhead
\bottomrule\noalign{}
\endlastfoot
Stop-and-Wait ARQ & 1 & 1 & The one frame & Individual; seq numbers 0/1 \\
\textbf{Go-Back-N} & 2ⁿ − 1 & 1 & Lost frame \textbf{and all after it} & Cumulative \\
\textbf{Selective Repeat} & 2ⁿ⁻¹ & 2ⁿ⁻¹ & Only lost frame & Individual / NAK \\
\end{longtable}

\endgroup

\textbf{Worked}: 1 Mbps link, frame 1000 bits, Tp = 20 ms. Tt = 1 ms, a = 20.
Stop-and-Wait η = 1/41 ≈ \textbf{2.4\%}. Window for 100\% = 1 + 40 = 41 → sequence bits = ⌈log₂ 41⌉ = \textbf{6} (GBN; 2⁶ − 1 = 63 ≥ 41).

\needspace{95pt}

\hypertarget{p2-09--64-hdlc--ppp}{%
\subsection{HDLC \& PPP}\label{p2-09--64-hdlc--ppp}}

\begin{itemize}
\tightlist
\item
  \textbf{HDLC}: bit-oriented; frames I (information), S (supervisory), U (unnumbered); modes NRM, ABM, ARM; flag 01111110.
\item
  \textbf{PPP}: byte-oriented, point-to-point (dial-up, DSL); LCP (link control), NCP, authentication \textbf{PAP} (plain-text) and \textbf{CHAP} (challenge-handshake, more secure). No flow control/sequence numbers by default.
\end{itemize}

\needspace{95pt}

\hypertarget{p2-09--7-medium-access-control}{%
\section{Medium Access Control}\label{p2-09--7-medium-access-control}}

\diagram{figures/p2-09-02}

\begin{listingblock}{62pt}
\begin{Verbatim}[fontsize=\fontsize{8.6}{10.4}\selectfont]
Pure ALOHA:     S = G·e^(−2G), max 1/(2e) = 0.184 at G = 0.5 ; vulnerable time = 2·Tt
Slotted ALOHA:  S = G·e^(−G),  max 1/e  = 0.368 at G = 1   ; vulnerable time = Tt
CSMA/CD:        minimum frame size: Tt ≥ 2·Tp  ⇒  L_min = 2 · Tp · B
                Efficiency = 1 / (1 + 6.44a)
\end{Verbatim}
\end{listingblock}

\textbf{Worked}: 10 Mbps Ethernet, max RTT 51.2 µs → L\_min = 10⁷ × 51.2 × 10⁻⁶ = \textbf{512 bits = 64 bytes}.

\begin{itemize}
\tightlist
\item
  CSMA/CD: on collision send \textbf{jam signal}, use \textbf{binary exponential backoff} (wait random k × slot, k ∈ {[}0, 2ⁿ − 1{]}; give up after 16 attempts).
\item
  CSMA/CA (802.11): \textbf{IFS} (DIFS, SIFS), \textbf{contention window}, \textbf{RTS/CTS} (solves \textbf{hidden terminal}), NAV, ACK.
\item
  Token ring (802.5), token bus (802.4).
\end{itemize}

\needspace{101pt}

\hypertarget{p2-09--ethernet-ieee-8023}{%
\subsection{Ethernet (IEEE 802.3)}\label{p2-09--ethernet-ieee-8023}}

\begin{listingblock}{51pt}
\begin{Verbatim}[fontsize=\fontsize{8.5}{10.3}\selectfont]
Frame: | Preamble 7 | SFD 1 | Dest 6 | Src 6 | Type/Len 2 | Data 46–1500 | CRC 4 |
Min frame = 64 B (without preamble) ; Max = 1518 B ; MAC address 48 bits (OUI 24 + 24)
Broadcast MAC: FF:FF:FF:FF:FF:FF ; multicast: LSB of first byte = 1
\end{Verbatim}
\end{listingblock}

\begingroup\small

\begin{longtable}[]{@{}
  >{\raggedright\arraybackslash}p{(\columnwidth - 4\tabcolsep) * \real{0.1943}}
  >{\raggedright\arraybackslash}p{(\columnwidth - 4\tabcolsep) * \real{0.1011}}
  >{\raggedright\arraybackslash}p{(\columnwidth - 4\tabcolsep) * \real{0.1661}}@{}}
\toprule\noalign{}
\begin{minipage}[b]{\linewidth}\raggedright
\textbf{Standard}
\end{minipage} & \begin{minipage}[b]{\linewidth}\raggedright
\textbf{Speed}
\end{minipage} & \begin{minipage}[b]{\linewidth}\raggedright
\textbf{Media}
\end{minipage} \\
\midrule\noalign{}
\endhead
\bottomrule\noalign{}
\endlastfoot
10Base5 / 10Base2 & 10 Mbps & Thick / thin coax \\
10Base-T & 10 Mbps & UTP \\
100Base-TX / FX & 100 Mbps & UTP / fibre \\
1000Base-T & 1 Gbps & UTP Cat5e \\
10GBase & 10 Gbps & Fibre/Cat6a \\
\end{longtable}

\endgroup

Devices: \textbf{Repeater/hub} (L1, one collision \& broadcast domain), \textbf{Bridge/switch} (L2, separates collision domains, one broadcast domain; learning, STP), \textbf{Router} (L3, separates broadcast domains), \textbf{Gateway} (all layers / protocol conversion). VLANs split broadcast domains on switches.

\needspace{166pt}

\hypertarget{p2-09--8-network-layer}{%
\section{Network Layer}\label{p2-09--8-network-layer}}

\needspace{130pt}

\hypertarget{p2-09--81-ipv4-addressing}{%
\subsection{IPv4 Addressing}\label{p2-09--81-ipv4-addressing}}

\begingroup\small

\begin{longtable}[]{@{}
  >{\raggedright\arraybackslash}p{(\columnwidth - 10\tabcolsep) * \real{0.0824}}
  >{\raggedright\arraybackslash}p{(\columnwidth - 10\tabcolsep) * \real{0.1276}}
  >{\raggedright\arraybackslash}p{(\columnwidth - 10\tabcolsep) * \real{0.1391}}
  >{\raggedright\arraybackslash}p{(\columnwidth - 10\tabcolsep) * \real{0.1635}}
  >{\raggedright\arraybackslash}p{(\columnwidth - 10\tabcolsep) * \real{0.1814}}
  >{\raggedright\arraybackslash}p{(\columnwidth - 10\tabcolsep) * \real{0.1262}}@{}}
\toprule\noalign{}
\begin{minipage}[b]{\linewidth}\raggedright
\textbf{Class}
\end{minipage} & \begin{minipage}[b]{\linewidth}\raggedright
\textbf{First bits}
\end{minipage} & \begin{minipage}[b]{\linewidth}\raggedright
\textbf{First octet}
\end{minipage} & \begin{minipage}[b]{\linewidth}\raggedright
\textbf{Default mask}
\end{minipage} & \begin{minipage}[b]{\linewidth}\raggedright
\textbf{Networks}
\end{minipage} & \begin{minipage}[b]{\linewidth}\raggedright
\textbf{Hosts/net}
\end{minipage} \\
\midrule\noalign{}
\endhead
\bottomrule\noalign{}
\endlastfoot
A & 0 & 0--127 & /8 255.0.0.0 & 2⁷ (126 usable) & 2²⁴ − 2 \\
B & 10 & 128--191 & /16 & 2¹⁴ & 2¹⁶ − 2 \\
C & 110 & 192--223 & /24 & 2²¹ & 254 \\
D & 1110 & 224--239 & --- & Multicast & --- \\
E & 1111 & 240--255 & --- & Reserved & --- \\
\end{longtable}

\endgroup

\begin{itemize}
\tightlist
\item
  Private ranges: \textbf{10.0.0.0/8, 172.16.0.0/12, 192.168.0.0/16}. Loopback: 127.0.0.0/8. APIPA 169.254/16.
\item
  \textbf{Subnetting worked}: 192.168.10.0/26
\end{itemize}

\begin{listingblock}{86pt}
\begin{Verbatim}[fontsize=\fontsize{9.0}{10.9}\selectfont]
Mask 255.255.255.192 → 4 subnets of 64 addresses (62 usable hosts each)
Subnet 1: 192.168.10.0   – .63   (broadcast .63)
Subnet 2: 192.168.10.64  – .127
Subnet 3: 192.168.10.128 – .191
Subnet 4: 192.168.10.192 – .255
Host 192.168.10.100/26 → network 192.168.10.64, broadcast 192.168.10.127
\end{Verbatim}
\end{listingblock}

\begin{itemize}
\tightlist
\item
  Usable hosts in /n = 2\^{}(32−n) − 2 (except /31, /32).
\item
  \textbf{CIDR} \& supernetting: 200.1.0.0/24 \ldots{} 200.1.3.0/24 aggregate to \textbf{200.1.0.0/22}. Longest prefix match in forwarding.
\item
  \textbf{NAT}: maps private addresses to public; PAT uses ports.
\end{itemize}

\needspace{124pt}

\hypertarget{p2-09--82-ipv4-header}{%
\subsection{IPv4 Header}\label{p2-09--82-ipv4-header}}

\begin{listingblock}{74pt}
\begin{Verbatim}[fontsize=\fontsize{8.8}{10.7}\selectfont]
| Ver 4 | HLEN 4 | DSCP/ToS 8 | Total length 16 |
| Identification 16 | Flags 3 (–, DF, MF) | Fragment offset 13 (units of 8 bytes) |
| TTL 8 | Protocol 8 (1 ICMP, 6 TCP, 17 UDP) | Header checksum 16 |
| Source IP 32 | Destination IP 32 | Options (0–40 B) |
Header 20–60 bytes (HLEN × 4) ; max datagram 65,535 bytes
\end{Verbatim}
\end{listingblock}

\textbf{Fragmentation worked}: 4000-byte datagram (20 B header) over MTU 1500 → data per fragment = 1480 (multiple of 8):
fragments carry 1480, 1480, 1020 bytes; offsets \textbf{0, 185, 370}; MF = 1, 1, 0.

\needspace{95pt}

\hypertarget{p2-09--83-ipv6}{%
\subsection{IPv6}\label{p2-09--83-ipv6}}

\begin{itemize}
\tightlist
\item
  128-bit address (hex, colon notation; \texttt{::} compresses zeros once), fixed \textbf{40-byte header}, no header checksum, \textbf{no fragmentation at routers} (source only), extension headers, flow label, built-in IPsec support, types: unicast, multicast, \textbf{anycast} (no broadcast).
\item
  Transition: dual stack, tunnelling, header translation.
\end{itemize}

\needspace{130pt}

\hypertarget{p2-09--84-routing}{%
\subsection{Routing}\label{p2-09--84-routing}}

\begingroup\small

\begin{longtable}[]{@{}
  >{\raggedright\arraybackslash}p{(\columnwidth - 6\tabcolsep) * \real{0.1155}}
  >{\raggedright\arraybackslash}p{(\columnwidth - 6\tabcolsep) * \real{0.3655}}
  >{\raggedright\arraybackslash}p{(\columnwidth - 6\tabcolsep) * \real{0.3351}}
  >{\raggedright\arraybackslash}p{(\columnwidth - 6\tabcolsep) * \real{0.1839}}@{}}
\toprule\noalign{}
\begin{minipage}[b]{\linewidth}\raggedright
\end{minipage} & \begin{minipage}[b]{\linewidth}\raggedright
\textbf{Distance Vector}
\end{minipage} & \begin{minipage}[b]{\linewidth}\raggedright
\textbf{Link State}
\end{minipage} & \begin{minipage}[b]{\linewidth}\raggedright
\textbf{Path Vector}
\end{minipage} \\
\midrule\noalign{}
\endhead
\bottomrule\noalign{}
\endlastfoot
Algorithm & \textbf{Bellman--Ford} & \textbf{Dijkstra} & --- \\
Knowledge & Neighbours\textquotesingle{} tables & Whole topology (LSPs flooded) & Paths (AS list) \\
Protocol & \textbf{RIP} & \textbf{OSPF} (also IS-IS) & \textbf{BGP} \\
Problem & \textbf{Count-to-infinity} (fix: split horizon, poison reverse, hold-down) & Heavy computation/\allowbreak{}flooding & Policy-based \\
Scope & Intra-AS (IGP) & Intra-AS (areas, backbone area 0) & Inter-AS (EGP) \\
\end{longtable}

\endgroup

\begin{itemize}
\tightlist
\item
  RIP: hop count metric, max \textbf{15} (16 = infinity), updates every \textbf{30 s}, runs over \textbf{UDP 520}.
\item
  OSPF: cost metric, runs directly over \textbf{IP (protocol 89)}, hierarchical areas.
\item
  BGP: runs over \textbf{TCP 179}; eBGP \& iBGP.
\item
  Multicast routing: DVMRP, MOSPF, PIM; \textbf{IGMP} for group membership.
\end{itemize}

\needspace{130pt}

\hypertarget{p2-09--85-support-protocols}{%
\subsection{Support Protocols}\label{p2-09--85-support-protocols}}

\begingroup\small

\begin{longtable}[]{@{}
  >{\raggedright\arraybackslash}p{(\columnwidth - 2\tabcolsep) * \real{0.2539}}
  >{\raggedright\arraybackslash}p{(\columnwidth - 2\tabcolsep) * \real{0.7461}}@{}}
\toprule\noalign{}
\begin{minipage}[b]{\linewidth}\raggedright
\textbf{Protocol}
\end{minipage} & \begin{minipage}[b]{\linewidth}\raggedright
\textbf{Purpose}
\end{minipage} \\
\midrule\noalign{}
\endhead
\bottomrule\noalign{}
\endlastfoot
\textbf{ARP} & IP → MAC (broadcast request, unicast reply) \\
RARP / BOOTP / \textbf{DHCP} & Obtain IP address (DHCP: DORA --- Discover, Offer, Request, Ack; UDP 67/68) \\
\textbf{ICMP} & Error reporting \& query: echo (ping), destination unreachable, time exceeded (traceroute), redirect, source quench \\
IGMP & Multicast group membership \\
\end{longtable}

\endgroup

\needspace{130pt}

\hypertarget{p2-09--9-transport-layer}{%
\section{Transport Layer}\label{p2-09--9-transport-layer}}

\begingroup\small

\begin{longtable}[]{@{}
  >{\raggedright\arraybackslash}p{(\columnwidth - 6\tabcolsep) * \real{0.2369}}
  >{\raggedright\arraybackslash}p{(\columnwidth - 6\tabcolsep) * \real{0.2268}}
  >{\raggedright\arraybackslash}p{(\columnwidth - 6\tabcolsep) * \real{0.2909}}
  >{\raggedright\arraybackslash}p{(\columnwidth - 6\tabcolsep) * \real{0.2454}}@{}}
\toprule\noalign{}
\begin{minipage}[b]{\linewidth}\raggedright
\textbf{Feature}
\end{minipage} & \begin{minipage}[b]{\linewidth}\raggedright
\textbf{TCP}
\end{minipage} & \begin{minipage}[b]{\linewidth}\raggedright
\textbf{UDP}
\end{minipage} & \begin{minipage}[b]{\linewidth}\raggedright
\textbf{SCTP}
\end{minipage} \\
\midrule\noalign{}
\endhead
\bottomrule\noalign{}
\endlastfoot
Connection & Connection-oriented & Connectionless & Connection-oriented (association) \\
Reliability & Reliable, ordered & Unreliable & Reliable, multi-stream \\
Unit & Byte stream & Message & Message (chunks) \\
Header & 20--60 B & \textbf{8 B} & 12 B common header \\
Flow/\allowbreak{}congestion control & Yes & No & Yes \\
Use & HTTP, FTP, SMTP & DNS, SNMP, DHCP, VoIP, TFTP & Signalling; multihoming \\
\end{longtable}

\endgroup

\needspace{95pt}

\hypertarget{p2-09--tcp-connection-management}{%
\subsection{TCP Connection Management}\label{p2-09--tcp-connection-management}}

\diagram{figures/p2-09-03}

\begin{itemize}
\tightlist
\item
  SYN and FIN each \textbf{consume one sequence number}. Sequence numbers count bytes.
\item
  TCP header flags: URG, ACK, PSH, RST, SYN, FIN. Window field 16 bits (window scaling option extends).
\item
  \textbf{Wrap-around time} = 2³² / bandwidth (bytes/s).
\item
  \textbf{Silly window syndrome}: Nagle\textquotesingle s algorithm (sender), Clark\textquotesingle s solution (receiver).
\item
  Timers: retransmission (RTO via Jacobson/Karn), persistence, keep-alive, TIME-WAIT.
\item
  RTT estimate: EstimatedRTT = (1 − α)·EstRTT + α·SampleRTT (α = 1/8).
\end{itemize}

\needspace{161pt}

\hypertarget{p2-09--tcp-congestion-control}{%
\subsection{TCP Congestion Control}\label{p2-09--tcp-congestion-control}}

\begin{listingblock}{111pt}
\begin{Verbatim}[fontsize=\fontsize{8.3}{10.1}\selectfont]
 cwnd
  │                 timeout → ssthresh = cwnd/2, cwnd = 1 (Tahoe/Reno)
  │            ╱╲   3 dup ACKs → ssthresh = cwnd/2, cwnd = ssthresh (Reno fast recovery)
  │          ╱    ╲
  │      ___╱       ╲_____/‾‾‾ (additive increase, +1 MSS per RTT)
  │     ╱ ← threshold
  │    ╱  slow start (exponential: doubles each RTT)
  │  ╱
  └─────────────────────────────────── RTT
\end{Verbatim}
\end{listingblock}

\textbf{Worked}: ssthresh = 16, starting cwnd = 1 MSS: 1, 2, 4, 8, 16 (slow start), then 17, 18, 19 \ldots{} (congestion avoidance). If timeout at cwnd = 20: ssthresh = 10, cwnd = 1.

\begin{itemize}
\tightlist
\item
  Effective window = min(cwnd, rwnd).
\item
  \textbf{QoS}: leaky bucket (constant output rate), \textbf{token bucket} (allows bursts: max burst S = C / (M − ρ)), IntServ (RSVP), DiffServ.
\end{itemize}

\needspace{101pt}

\hypertarget{p2-09--well-known-ports}{%
\subsection{Well-known Ports}\label{p2-09--well-known-ports}}

\begin{listingblock}{51pt}
\begin{Verbatim}[fontsize=\fontsize{8.5}{10.3}\selectfont]
FTP 20 (data) / 21 (control)   SSH 22   TELNET 23   SMTP 25   DNS 53   DHCP 67/68
TFTP 69   HTTP 80   POP3 110   NTP 123   IMAP 143   SNMP 161/162   BGP 179   HTTPS 443
Ports: well-known 0–1023, registered 1024–49151, dynamic 49152–65535
\end{Verbatim}
\end{listingblock}

\needspace{95pt}

\hypertarget{p2-09--10-application-layer}{%
\section{Application Layer}\label{p2-09--10-application-layer}}

\begin{itemize}
\tightlist
\item
  \textbf{DNS}: hierarchical (root → TLD → authoritative); records \textbf{A} (IPv4), \textbf{AAAA} (IPv6), \textbf{MX} (mail), \textbf{CNAME} (alias), \textbf{NS}, \textbf{PTR} (reverse), SOA; recursive vs iterative resolution; UDP 53 (TCP for zone transfers). \textbf{DDNS} updates records dynamically.
\item
  \textbf{HTTP}: stateless; non-persistent (HTTP/1.0; 2 RTT per object) vs persistent (HTTP/1.1, pipelining); cookies for state; methods GET, POST, PUT, DELETE, HEAD; status 1xx info, 2xx success (200 OK), 3xx redirect (301), 4xx client error (404), 5xx server error (500). HTTP/2 multiplexing, HTTP/3 over QUIC (UDP).
\item
  \textbf{FTP}: two connections --- control (21, persistent, out-of-band) \& data (20, per file).
\item
  \textbf{Email}: SMTP (push), POP3/IMAP (pull), MIME for attachments. See \protect\hyperlink{p1-08}{Paper 1 ICT}.
\item
  \textbf{TELNET} (insecure remote login; NVT), \textbf{SSH} (secure).
\item
  \textbf{SNMP}: manager--agent, \textbf{MIB}, \textbf{SMI}; GET, SET, TRAP; UDP 161/162.
\item
  \textbf{WWW}: URL = protocol://host:port/path; web caching (proxy).
\item
  \textbf{Bluetooth} (IEEE 802.15.1): piconet (1 primary + up to \textbf{7 active} secondaries), scatternet; 2.4 GHz FHSS.
\end{itemize}

\needspace{197pt}

\hypertarget{p2-09--11-network-security}{%
\section{Network Security}\label{p2-09--11-network-security}}

\needspace{161pt}

\hypertarget{p2-09--111-goals--attacks}{%
\subsection{Goals \& Attacks}\label{p2-09--111-goals--attacks}}

CIA: \textbf{Confidentiality, Integrity, Availability} + authentication, non-repudiation.
Passive attacks (eavesdropping, traffic analysis) vs active (masquerade, replay, modification, DoS).
Malware: virus, worm, Trojan, ransomware, spyware, rootkit, botnet, logic bomb, keylogger.

\needspace{95pt}

\hypertarget{p2-09--112-cryptography}{%
\subsection{Cryptography}\label{p2-09--112-cryptography}}

\diagram{figures/p2-09-04}

\begingroup\small

\begin{longtable}[]{@{}
  >{\raggedright\arraybackslash}p{(\columnwidth - 4\tabcolsep) * \real{0.2609}}
  >{\raggedright\arraybackslash}p{(\columnwidth - 4\tabcolsep) * \real{0.2897}}
  >{\raggedright\arraybackslash}p{(\columnwidth - 4\tabcolsep) * \real{0.4494}}@{}}
\toprule\noalign{}
\begin{minipage}[b]{\linewidth}\raggedright
\textbf{Algorithm}
\end{minipage} & \begin{minipage}[b]{\linewidth}\raggedright
\textbf{Type}
\end{minipage} & \begin{minipage}[b]{\linewidth}\raggedright
\textbf{Key facts}
\end{minipage} \\
\midrule\noalign{}
\endhead
\bottomrule\noalign{}
\endlastfoot
Caesar & Substitution (mono-alphabetic) & Shift by k (k = 3) \\
Vigenère & Poly-alphabetic & Keyword shifts \\
Playfair & Digraph substitution & 5×5 matrix \\
Rail fence, columnar & \textbf{Transposition} & Rearranges letters \\
\textbf{DES} & Symmetric block (Feistel) & 64-bit block, \textbf{56-bit key}, \textbf{16 rounds} \\
3DES & Symmetric & EDE, 112/168-bit key \\
\textbf{AES (Rijndael)} & Symmetric block (SPN) & 128-bit block; keys 128/192/256 → \textbf{10/12/14 rounds} \\
IDEA & Symmetric & 64-bit block, 128-bit key \\
RC4 & Stream cipher & WEP \\
\textbf{RSA} & Asymmetric & Factoring problem \\
\textbf{Diffie--Hellman} & Key exchange & Discrete log; vulnerable to man-in-the-middle \\
ECC & Asymmetric & Smaller keys for same security \\
MD5 / SHA-1 / SHA-256 & Hash & 128 / 160 / 256-bit digests \\
\end{longtable}

\endgroup

\begin{itemize}
\tightlist
\item
  Number of keys for n users: symmetric \textbf{n(n − 1)/2}, asymmetric \textbf{2n}.
\item
  \textbf{RSA worked}: p = 3, q = 11 → n = 33, φ = 20; choose e = 3 (gcd(3,20)=1); d = 7 (3×7 = 21 ≡ 1 mod 20). Encrypt M = 4: C = 4³ mod 33 = 64 mod 33 = \textbf{31}. Decrypt: 31⁷ mod 33 = \textbf{4}.
\item
  \textbf{Diffie--Hellman worked}: p = 23, g = 5, a = 6, b = 15 → A = 5⁶ mod 23 = 8, B = 5¹⁵ mod 23 = 19; shared key = 19⁶ mod 23 = 8¹⁵ mod 23 = \textbf{2}.
\end{itemize}

\needspace{101pt}

\hypertarget{p2-09--113-digital-signature--pki}{%
\subsection{Digital Signature \& PKI}\label{p2-09--113-digital-signature--pki}}

\begin{listingblock}{51pt}
\begin{Verbatim}[fontsize=\fontsize{8.3}{10.1}\selectfont]
Sign:   signature = Encrypt(hash(M), sender's PRIVATE key)
Verify: Decrypt(signature, sender's PUBLIC key) == hash(M)
Confidentiality: encrypt with receiver's PUBLIC key; decrypt with receiver's PRIVATE key
\end{Verbatim}
\end{listingblock}

Provides authentication, integrity, \textbf{non-repudiation}. Certificates (X.509) issued by \textbf{CA}; MAC/HMAC for message integrity with shared key.

\begin{itemize}
\tightlist
\item
  \textbf{Steganography}: hiding the existence of a message (e.g. in image LSBs), vs cryptography (hides meaning).
\end{itemize}

\needspace{95pt}

\hypertarget{p2-09--114-security-protocols--devices}{%
\subsection{Security Protocols \& Devices}\label{p2-09--114-security-protocols--devices}}

\begin{itemize}
\tightlist
\item
  \textbf{IPsec} (network layer): AH (authentication), ESP (encryption); transport vs tunnel mode; IKE.
\item
  \textbf{SSL/TLS} (transport layer): handshake, record protocol; HTTPS. \textbf{PGP} \& S/MIME (email).
\item
  \textbf{VPN}: secure tunnel over public network (IPsec, SSL VPN, PPTP, L2TP).
\item
  \textbf{Firewalls}: packet filter (network/transport headers --- stateless), stateful inspection, \textbf{proxy / application-level gateway}, circuit-level gateway. DMZ. \textbf{IDS/IPS} (signature vs anomaly).
\item
  Kerberos (tickets, KDC: AS + TGS).
\end{itemize}

\needspace{131pt}

\hypertarget{p2-09--12-mobile--wireless}{%
\section{Mobile \& Wireless}\label{p2-09--12-mobile--wireless}}

\needspace{95pt}

\hypertarget{p2-09--121-cellular-concepts}{%
\subsection{Cellular Concepts}\label{p2-09--121-cellular-concepts}}

\begin{itemize}
\tightlist
\item
  Frequency reuse; cluster size N = i² + ij + j² (1, 3, 4, 7, 12\ldots); co-channel reuse distance D = R√(3N).
\item
  Handoff (hard --- GSM, soft --- CDMA). Generations: 1G (analog AMPS), 2G (GSM/CDMA digital voice, SMS), 2.5G (\textbf{GPRS}), 2.75G (EDGE), 3G (UMTS/WCDMA, CDMA2000), 4G (LTE, all-IP, OFDMA), 5G (mmWave, massive MIMO, URLLC, eMBB, mMTC).
\end{itemize}

\needspace{95pt}

\hypertarget{p2-09--122-gsm-architecture}{%
\subsection{GSM Architecture}\label{p2-09--122-gsm-architecture}}

\diagram{figures/p2-09-05}

\begin{itemize}
\tightlist
\item
  GSM: TDMA + FDMA; 900/1800 MHz; 200 kHz carriers, 8 time slots; \textbf{HLR} (permanent subscriber data), \textbf{VLR} (temporary for visitors), \textbf{AuC} (authentication keys), \textbf{EIR} (IMEI database). IMSI in SIM.
\item
  \textbf{GPRS}: packet-switched over GSM; SGSN \& GGSN nodes. \textbf{SMS}: 160 characters, via SMSC, uses signalling channels.
\item
  \textbf{CDMA} (IS-95): spread spectrum, unique codes, soft handoff, RAKE receiver.
\end{itemize}

\needspace{95pt}

\hypertarget{p2-09--123-mobile-ip--others}{%
\subsection{Mobile IP \& Others}\label{p2-09--123-mobile-ip--others}}

\begin{itemize}
\tightlist
\item
  \textbf{Mobile IP}: home agent, foreign agent, \textbf{care-of address} (CoA), registration, \textbf{tunnelling}, triangle routing problem (route optimisation).
\item
  Mobile computing middleware \& gateways (WAP; WAP gateway translates between WAP \& HTTP).
\item
  \textbf{MANET}: infrastructure-less, self-configuring; routing --- proactive (\textbf{DSDV}, OLSR), reactive (\textbf{AODV}, \textbf{DSR}), hybrid (ZRP).
\item
  \textbf{Wireless LAN (IEEE 802.11)}: infrastructure (BSS, ESS via APs) vs ad hoc (IBSS); 802.11a/b/g/n/ac/ax (Wi-Fi 6). Security WEP → WPA → WPA2 (AES-CCMP) → WPA3.
\item
  \textbf{Satellites}: GEO (35,786 km; \textasciitilde270 ms one-way delay; 3 satellites cover earth), MEO (GPS \textasciitilde20,200 km), LEO (500--2000 km; Starlink, Iridium). Uplink frequency \textgreater{} downlink.
\item
  \textbf{Wireless geolocation}: GPS (trilateration, ≥ 4 satellites for 3-D fix), cell-ID, AoA, ToA, TDoA, RSSI. India: \textbf{NavIC (IRNSS, 7 satellites)}.
\end{itemize}

\needspace{158pt}

\hypertarget{p2-09--13-cloud-computing--iot}{%
\section{Cloud Computing \& IoT}\label{p2-09--13-cloud-computing--iot}}

\begin{listingblock}{108pt}
\begin{Verbatim}[fontsize=\fontsize{9.0}{10.9}\selectfont]
   ┌─────────────────────────────────────────────────────────────┐
   │ SaaS  (Gmail, Salesforce, Office 365)    — user uses app    │
   ├─────────────────────────────────────────────────────────────┤
   │ PaaS  (Google App Engine, Heroku, Azure App Service) — dev  │
   ├─────────────────────────────────────────────────────────────┤
   │ IaaS  (AWS EC2, Google Compute Engine, Azure VMs) — admin   │
   └─────────────────────────────────────────────────────────────┘
   Customer manages more as you go down ; provider manages more as you go up
\end{Verbatim}
\end{listingblock}

\begin{itemize}
\tightlist
\item
  \textbf{NIST essential characteristics (5)}: on-demand self-service, broad network access, resource pooling, rapid elasticity, measured service.
\item
  Deployment: \textbf{public, private, community, hybrid}.
\item
  \textbf{Virtualisation}: hypervisors (see \protect\hyperlink{p2-05--11-virtual-machines}{OS unit}); virtual servers; live migration.
\item
  Cloud storage (object: S3; block: EBS; file), database-as-a-service; multi-tenancy; resource management (load balancing, auto-scaling).
\item
  \textbf{SLA} (Service Level Agreement): availability (e.g. 99.9\% → \textasciitilde8.76 h downtime/year), response time, penalties.
\item
  \textbf{IoT}: things with sensors + connectivity. Architecture (3-layer): \textbf{Perception (sensing) → Network → Application} (5-layer adds processing/middleware \& business). Protocols: \textbf{MQTT} (publish--subscribe over TCP), \textbf{CoAP} (REST over UDP), 6LoWPAN, Zigbee (802.15.4), BLE, LoRaWAN, NB-IoT, RFID/NFC. Edge/fog computing reduces latency.
\end{itemize}

\needspace{196pt}

\hypertarget{p2-09--14-deeper-dive--vlsm-routing-tables-tcp-numbers-hamming--ciphers}{%
\section{Deeper Dive --- VLSM, Routing Tables, TCP Numbers, Hamming \& Ciphers}\label{p2-09--14-deeper-dive--vlsm-routing-tables-tcp-numbers-hamming--ciphers}}

\needspace{160pt}

\hypertarget{p2-09--141-vlsm--worked}{%
\subsection{VLSM --- Worked}\label{p2-09--141-vlsm--worked}}

Split \textbf{192.168.1.0/24} for LANs needing 100, 50, 25 and 10 hosts (allocate largest first):

\begingroup\small

\begin{longtable}[]{@{}
  >{\raggedright\arraybackslash}p{(\columnwidth - 10\tabcolsep) * \real{0.0774}}
  >{\raggedright\arraybackslash}p{(\columnwidth - 10\tabcolsep) * \real{0.1592}}
  >{\raggedright\arraybackslash}p{(\columnwidth - 10\tabcolsep) * \real{0.1872}}
  >{\raggedright\arraybackslash}p{(\columnwidth - 10\tabcolsep) * \real{0.1782}}
  >{\raggedright\arraybackslash}p{(\columnwidth - 10\tabcolsep) * \real{0.1474}}
  >{\raggedright\arraybackslash}p{(\columnwidth - 10\tabcolsep) * \real{0.1262}}@{}}
\toprule\noalign{}
\begin{minipage}[b]{\linewidth}\raggedright
\textbf{LAN}
\end{minipage} & \begin{minipage}[b]{\linewidth}\raggedright
\textbf{Hosts needed}
\end{minipage} & \begin{minipage}[b]{\linewidth}\raggedright
\textbf{Prefix (usable)}
\end{minipage} & \begin{minipage}[b]{\linewidth}\raggedright
\textbf{Subnet}
\end{minipage} & \begin{minipage}[b]{\linewidth}\raggedright
\textbf{Range}
\end{minipage} & \begin{minipage}[b]{\linewidth}\raggedright
\textbf{Broadcast}
\end{minipage} \\
\midrule\noalign{}
\endhead
\bottomrule\noalign{}
\endlastfoot
A & 100 & /25 (126) & 192.168.1.0 & .1 -- .126 & .127 \\
B & 50 & /26 (62) & 192.168.1.128 & .129 -- .190 & .191 \\
C & 25 & /27 (30) & 192.168.1.192 & .193 -- .222 & .223 \\
D & 10 & /28 (14) & 192.168.1.224 & .225 -- .238 & .239 \\
Free & --- & /28 & 192.168.1.240 & --- & --- \\
\end{longtable}

\endgroup

\needspace{95pt}

\hypertarget{p2-09--142-distance-vector-update--worked}{%
\subsection{Distance-Vector Update --- Worked}\label{p2-09--142-distance-vector-update--worked}}

\diagram{figures/p2-09-06}

\begingroup\small

\begin{longtable}[]{@{}
  >{\raggedright\arraybackslash}p{(\columnwidth - 4\tabcolsep) * \real{0.1099}}
  >{\raggedright\arraybackslash}p{(\columnwidth - 4\tabcolsep) * \real{0.1074}}
  >{\raggedright\arraybackslash}p{(\columnwidth - 4\tabcolsep) * \real{0.4833}}@{}}
\toprule\noalign{}
\begin{minipage}[b]{\linewidth}\raggedright
\textbf{A\textquotesingle s table}
\end{minipage} & \begin{minipage}[b]{\linewidth}\raggedright
\textbf{Initially}
\end{minipage} & \begin{minipage}[b]{\linewidth}\raggedright
\textbf{After receiving B\textquotesingle s vector (B: A 1, B 0, C 2)}
\end{minipage} \\
\midrule\noalign{}
\endhead
\bottomrule\noalign{}
\endlastfoot
to A & 0 & 0 \\
to B & 1 (direct) & 1 (direct) \\
to C & 5 (direct) & min(5, 1 + 2) = \textbf{3 via B} \\
\end{longtable}

\endgroup

Bellman--Ford equation: Dₓ(y) = minᵥ \{ c(x, v) + Dᵥ(y) \}.

\textbf{Count-to-infinity}: suppose link B--C fails. B may then believe A\textquotesingle s stale advertisement "C at cost 3" (which was itself via B) and set C = 1 + 3 = 4 via A; A then updates to 1 + 4 = 5 via B, and the two keep raising each other\textquotesingle s cost one step at a time. Here the climb stops once it exceeds A\textquotesingle s direct link (cost 5), but with no alternative path it would continue until "infinity" (16 in RIP). \textbf{Split horizon} (never advertise a route back to the neighbour it was learned from) and \textbf{poison reverse} (advertise it as ∞) cure this two-node loop.

\needspace{95pt}

\hypertarget{p2-09--143-tcp-sequence-and-acknowledgement-numbers}{%
\subsection{TCP Sequence and Acknowledgement Numbers}\label{p2-09--143-tcp-sequence-and-acknowledgement-numbers}}

\diagram{figures/p2-09-07}

Next expected byte = last seq + length; SYN and FIN consume one number each.

\textbf{Go-Back-N count}: window 4, frames 0--6, frame 2 lost once (no other losses). Assume frames 0--4 have been sent when frame 2\textquotesingle s timer expires (ACKs for 0 and 1 slid the window to 2--5).
Sent 0, 1, 2 (lost), 3, 4 → on timeout resend 2, 3, 4, then send 5, 6. Total transmissions = 5 + 3 + 2 = \textbf{10}. Selective Repeat resends only frame 2: 7 + 1 = \textbf{8}.

\needspace{220pt}

\hypertarget{p2-09--144-hamming-7-4--encode-and-correct}{%
\subsection{Hamming (7, 4) --- Encode and Correct}\label{p2-09--144-hamming-7-4--encode-and-correct}}

Data 1011 → positions: 1 p1, 2 p2, 3 d1, 4 p4, 5 d2, 6 d3, 7 d4 (even parity).

\begin{listingblock}{140pt}
\begin{Verbatim}[fontsize=\fontsize{9.0}{10.9}\selectfont]
d1 d2 d3 d4 = 1 0 1 1
p1 covers 1,3,5,7 → d1 d2 d4 = 1 0 1 → p1 = 0
p2 covers 2,3,6,7 → d1 d3 d4 = 1 1 1 → p2 = 1
p4 covers 4,5,6,7 → d2 d3 d4 = 0 1 1 → p4 = 0
Codeword (positions 1..7) = 0 1 1 0 0 1 1

Received 0 1 1 0 0 0 1 (bit 6 flipped):
c1 = b1 b3 b5 b7 = 0 1 0 1 → 0
c2 = b2 b3 b6 b7 = 1 1 0 1 → 1
c4 = b4 b5 b6 b7 = 0 0 0 1 → 1
Syndrome c4 c2 c1 = 110₂ = 6 → flip bit 6 → corrected
\end{Verbatim}
\end{listingblock}

\needspace{95pt}

\hypertarget{p2-09--145-delay-calculations}{%
\subsection{Delay Calculations}\label{p2-09--145-delay-calculations}}

\begin{itemize}
\tightlist
\item
  Store-and-forward: a 1000-byte packet over 3 links of 1 Mbps (propagation ignored) → 3 × 8 ms = \textbf{24 ms}.
\item
  Message of 3 packets over the same path (pipelining): (3 + 3 − 1) × 8 = \textbf{40 ms}.
\item
  Circuit switching instead: setup time + message/bandwidth + propagation.
\end{itemize}

\needspace{130pt}

\hypertarget{p2-09--146-ipv6-address-compression}{%
\subsection{IPv6 Address Compression}\label{p2-09--146-ipv6-address-compression}}

\begingroup\small

\begin{longtable}[]{@{}
  >{\raggedright\arraybackslash}p{(\columnwidth - 2\tabcolsep) * \real{0.4344}}
  >{\raggedright\arraybackslash}p{(\columnwidth - 2\tabcolsep) * \real{0.2478}}@{}}
\toprule\noalign{}
\begin{minipage}[b]{\linewidth}\raggedright
\textbf{Full}
\end{minipage} & \begin{minipage}[b]{\linewidth}\raggedright
\textbf{Compressed}
\end{minipage} \\
\midrule\noalign{}
\endhead
\bottomrule\noalign{}
\endlastfoot
2001:0db8:0000:0000:0000:ff00:0042:8329 & 2001:db8::ff00:42:8329 \\
fe80:0000:0000:0000:0202:b3ff:fe1e:8329 & fe80::202:b3ff:fe1e:8329 \\
0000:0000:0000:0000:0000:0000:0000:0001 & ::1 (loopback) \\
\end{longtable}

\endgroup

Rules: drop leading zeros in each group; replace \textbf{one} longest run of all-zero groups with \texttt{::}.

\needspace{95pt}

\hypertarget{p2-09--147-dns-resolution}{%
\subsection{DNS Resolution}\label{p2-09--147-dns-resolution}}

\diagram{figures/p2-09-08}

\needspace{95pt}

\hypertarget{p2-09--148-classical-ciphers--worked}{%
\subsection{Classical Ciphers --- Worked}\label{p2-09--148-classical-ciphers--worked}}

\begin{itemize}
\tightlist
\item
  \textbf{Caesar (k = 3)}: NET → \textbf{QHW}.
\item
  \textbf{Vigenère}, key LEMON: ATTACK → A+L = L, T+E = X, T+M = F, A+O = O, C+N = P, K+L = V → \textbf{LXFOPV}.
\item
  \textbf{Rail fence (2 rails)}: HELLOWORLD → rails HLOOL / ELWRD → \textbf{HLOOLELWRD}.
\item
  \textbf{Columnar transposition} rearranges letters by a key order; \textbf{substitution} replaces letters. Product ciphers (DES, AES) combine both (confusion + diffusion --- Shannon).
\end{itemize}

\needspace{136pt}

\hypertarget{p2-09--149-http-exchange-what-a-request-looks-like}{%
\subsection{HTTP Exchange (what a request looks like)}\label{p2-09--149-http-exchange-what-a-request-looks-like}}

\begin{listingblock}{86pt}
\begin{Verbatim}[fontsize=\fontsize{9.0}{10.9}\selectfont]
GET /index.html HTTP/1.1                 HTTP/1.1 200 OK
Host: www.ugcnet.example                 Content-Type: text/html
User-Agent: Mozilla/5.0                  Content-Length: 3421
Accept: text/html                        Set-Cookie: sid=abc123
Connection: keep-alive
                                         <html> … </html>
\end{Verbatim}
\end{listingblock}

\needspace{125pt}

\hypertarget{p2-09--previous-year-questions-pyq-pattern}{%
\section{Previous Year Questions (PYQ pattern)}\label{p2-09--previous-year-questions-pyq-pattern}}

\textbf{Models \& topologies}

\begin{enumerate}
\def\labelenumi{\arabic{enumi}.}
\tightlist
\item
  Which OSI layer handles encryption \& compression? \textbf{Presentation}
\item
  Process-to-process delivery is the job of: \textbf{Transport layer}
\item
  Dialog control \& synchronisation: \textbf{Session layer}
\item
  Number of links in a fully connected mesh of 10 nodes: \textbf{45}
\item
  Router operates at: \textbf{Network layer}; switch at: \textbf{Data link layer}
\item
  Number of collision domains with a 24-port switch: \textbf{24}; broadcast domains: \textbf{1}
\end{enumerate}

\textbf{Physical layer}

\begin{enumerate}
\def\labelenumi{\arabic{enumi}.}
\setcounter{enumi}{6}
\tightlist
\item
  Channel B = 4 kHz, SNR = 63: C = 4000 × log₂ 64 = \textbf{24 kbps}
\item
  Noiseless 3 kHz channel with 8 levels: 2 × 3000 × 3 = \textbf{18 kbps}
\item
  Manchester encoding needs bandwidth: \textbf{twice that of NRZ} (it is self-clocking)
\item
  Sampling a 4 kHz voice signal per Nyquist: \textbf{8000 samples/s}
\item
  16-QAM carries how many bits per symbol? \textbf{4}
\item
  T1 line data rate: \textbf{1.544 Mbps}
\item
  Medium immune to electromagnetic interference: \textbf{optical fibre}
\end{enumerate}

\textbf{Data link}

\begin{enumerate}
\def\labelenumi{\arabic{enumi}.}
\setcounter{enumi}{13}
\tightlist
\item
  Bit stuffing after five consecutive 1s inserts: \textbf{a 0}
\item
  CRC with generator of degree 4 appends how many bits? \textbf{4}
\item
  Min Hamming distance to correct 2-bit errors: \textbf{5}
\item
  Go-Back-N with 3-bit sequence numbers: max sender window \textbf{7}; Selective Repeat: \textbf{4}
\item
  Stop-and-wait, Tt = 1 ms, Tp = 2 ms: η = 1/(1 + 4) = \textbf{20\%}
\item
  Window size for full utilisation with a = 10: \textbf{21}
\item
  Which protocol uses cumulative ACKs and discards out-of-order frames? \textbf{Go-Back-N}
\item
  PPP authentication protocol using a challenge: \textbf{CHAP}
\end{enumerate}

\textbf{MAC}

\begin{enumerate}
\def\labelenumi{\arabic{enumi}.}
\setcounter{enumi}{21}
\tightlist
\item
  Max throughput of slotted ALOHA: \textbf{36.8\%}; pure ALOHA: \textbf{18.4\%}
\item
  CSMA/CD: 1 Gbps, 2 km cable, signal speed 2×10⁸ m/s → Tp = 10 µs → L\_min = 2 × 10 µs × 10⁹ = \textbf{20,000 bits}
\item
  Hidden terminal problem is solved by: \textbf{RTS/CTS (CSMA/CA)}
\item
  Minimum Ethernet frame size: \textbf{64 bytes}
\item
  Binary exponential backoff is used in: \textbf{CSMA/CD Ethernet}
\end{enumerate}

\textbf{Network layer}

\begin{enumerate}
\def\labelenumi{\arabic{enumi}.}
\setcounter{enumi}{26}
\tightlist
\item
  Class of 191.10.20.1: \textbf{Class B}
\item
  Usable hosts in a /27 subnet: \textbf{30}
\item
  Network address of 172.16.45.14/20: 45 = 0010 1101 → /20 keeps 0010 → \textbf{172.16.32.0}
\item
  Broadcast address of 10.1.1.130/25: \textbf{10.1.1.255}
\item
  Subnet mask for 500 hosts per subnet: \textbf{/23 (255.255.254.0)}
\item
  IPv4 header field preventing infinite looping: \textbf{TTL}
\item
  Fragment offset is measured in units of: \textbf{8 bytes}
\item
  IPv6 address size: \textbf{128 bits}; header: \textbf{40 bytes}
\item
  RIP uses: \textbf{distance vector, hop count (max 15)}
\item
  OSPF uses: \textbf{link state (Dijkstra)}
\item
  BGP is a: \textbf{path vector, inter-domain protocol over TCP 179}
\item
  Count-to-infinity occurs in: \textbf{distance vector routing}
\item
  Which protocol maps IP to MAC? \textbf{ARP}
\item
  \texttt{ping} and \texttt{traceroute} use: \textbf{ICMP}
\item
  DHCP message sequence: \textbf{Discover, Offer, Request, Acknowledge}
\end{enumerate}

\textbf{Transport \& application}

\begin{enumerate}
\def\labelenumi{\arabic{enumi}.}
\setcounter{enumi}{41}
\tightlist
\item
  UDP header size: \textbf{8 bytes}
\item
  TCP 3-way handshake segments: \textbf{SYN, SYN-ACK, ACK}
\item
  After a timeout with cwnd = 32, new ssthresh: \textbf{16}, cwnd: \textbf{1}
\item
  Port of SMTP \textbf{25}, DNS \textbf{53}, HTTP \textbf{80}, HTTPS \textbf{443}, FTP control \textbf{21}
\item
  DNS record for mail server: \textbf{MX}; for IPv6: \textbf{AAAA}
\item
  FTP uses \_\_ connections: \textbf{two (control + data)}; control channel is \textbf{out-of-band}
\item
  HTTP is: \textbf{stateless} (state maintained via cookies)
\item
  SNMP database of managed objects: \textbf{MIB}
\item
  Token bucket: capacity C = 1 MB, token rate ρ = 2 MBps, max rate M = 10 MBps → max burst time = C/(M − ρ) = \textbf{0.125 s}
\end{enumerate}

\textbf{Security}

\begin{enumerate}
\def\labelenumi{\arabic{enumi}.}
\setcounter{enumi}{50}
\tightlist
\item
  DES key length: \textbf{56 bits}; rounds: \textbf{16}
\item
  AES-128 rounds: \textbf{10}
\item
  RSA with p = 5, q = 11, e = 3 → φ = 40, d = \textbf{27}
\item
  In digital signatures, the sender signs with: \textbf{own private key}
\item
  Number of keys for 100 users with symmetric crypto: \textbf{4950}; asymmetric: \textbf{200}
\item
  Diffie--Hellman is vulnerable to: \textbf{man-in-the-middle attack}
\item
  Hiding a message inside an image: \textbf{steganography}
\item
  IPsec operates at: \textbf{network layer}; SSL/TLS at: \textbf{transport layer}
\item
  A firewall that examines application-layer data: \textbf{proxy (application-level gateway)}
\end{enumerate}

\textbf{Mobile, cloud, IoT}

\begin{enumerate}
\def\labelenumi{\arabic{enumi}.}
\setcounter{enumi}{59}
\tightlist
\item
  GSM database holding permanent subscriber information: \textbf{HLR}; temporary visitor data: \textbf{VLR}
\item
  GPRS is a: \textbf{packet-switched 2.5G service}
\item
  In Mobile IP, the address used at the foreign network: \textbf{care-of address}
\item
  AODV is a: \textbf{reactive (on-demand) MANET routing protocol}
\item
  GEO satellite altitude: \textbf{≈ 35,786 km}
\item
  Google App Engine is an example of: \textbf{PaaS}; AWS EC2: \textbf{IaaS}
\item
  Rapid elasticity is a characteristic of: \textbf{cloud computing (NIST)}
\item
  MQTT follows: \textbf{publish--subscribe model}
\item
  Bluetooth piconet can have at most \_\_ active secondaries: \textbf{7}
\end{enumerate}

\textbf{More practice questions}

\begin{enumerate}
\def\labelenumi{\arabic{enumi}.}
\setcounter{enumi}{68}
\tightlist
\item
  Using VLSM on 192.168.1.0/24, the prefix for a LAN with 50 hosts: \textbf{/26}
\item
  Broadcast address of 192.168.1.192/27: \textbf{192.168.1.223}
\item
  Split horizon is a remedy for: \textbf{count-to-infinity}
\item
  Client ISN 1000; after the handshake it sends 200 bytes. Sequence number of the next segment: \textbf{1201}
\item
  In the GBN scenario of §14.3 (frames 0--4 sent before the timeout), total transmissions: \textbf{10}; with Selective Repeat: \textbf{8}
\item
  Hamming (7, 4) codeword for data 1011 (even parity): \textbf{0110011}
\item
  Syndrome 101 in Hamming (7, 4) means the error is in bit: \textbf{5}
\item
  Compressed form of 2001:0db8:0000:0000:0000:0000:0000:0001: \textbf{2001:db8::1}
\item
  \texttt{::} may appear in an IPv6 address: \textbf{only once}
\item
  Vigenère encryption of "ATTACK" with key "LEMON": \textbf{LXFOPV}
\item
  Caesar cipher with key 3 encrypts "NET" as: \textbf{QHW}
\item
  Store-and-forward delay of a 1000-byte packet over 3 links of 1 Mbps: \textbf{24 ms}
\end{enumerate}

\begin{revbox}

\begin{itemize}
\tightlist
\item
  Shannon C = B log₂(1+SNR) · Nyquist 2B log₂L
\item
  η\_SW = 1/(1+2a) · GBN W ≤ 2ⁿ−1 · SR W ≤ 2ⁿ⁻¹ · L\_min = 2·Tp·B
\item
  ALOHA 18.4\% / 36.8\% · Ethernet min 64 B
\item
  /n hosts = 2\^{}(32−n) − 2 · RIP DV/15 hops/UDP 520 · OSPF LS/IP 89 · BGP PV/TCP 179
\item
  TCP: slow start exponential → AIMD; timeout → cwnd 1
\item
  DES 56/16 · AES 10/12/14 · sign with private, encrypt with receiver\textquotesingle s public
\end{itemize}

\end{revbox}
